\documentclass[11pt,a4paper]{article}
\usepackage[utf8]{inputenc}
\usepackage[vietnamese]{babel}
\usepackage{amsmath,amssymb}
\usepackage{geometry}
\geometry{a4paper, total={180mm,260mm}, left=15mm, right=15mm, top=18mm, bottom=18mm}
\usepackage{tikz}
\usetikzlibrary{shapes.geometric, arrows.meta, positioning, calc}
\usepackage{tcolorbox}
\tcbuselibrary{skins,breakable}
\usepackage{xcolor}
\usepackage{multicol}
\usepackage{array}
\usepackage{fontawesome}
\usepackage{fancyhdr}

\pagestyle{fancy}
\fancyhf{}
\renewcommand{\headrulewidth}{0.5pt}
\renewcommand{\footrulewidth}{0.5pt}
\fancyhead[L]{\small\bfseries\color{teal}\faBook\ SỔ TAY KIẾN THỨC TOÁN 3}
\fancyhead[R]{\small\bfseries\color{teal}Lớp Toán Thầy Lĩnh -- loptoanthaylinh.vn}
\fancyfoot[C]{\small\color{gray}-- Sổ tay kiến thức: Bài 66. Xem đồng hồ. Tháng -- năm --}

\definecolor{primary}{HTML}{005F73}
\definecolor{secondary}{HTML}{0A9396}
\definecolor{accent}{HTML}{EE9B00}
\definecolor{danger}{HTML}{AE2012}
\definecolor{boxbg}{HTML}{F4F6F8}

\begin{document}

% --- BANNER TIÊU ĐỀ ---
\begin{tcolorbox}[
  enhanced,
  colback=white,
  colframe=primary,
  arc=3mm,
  boxrule=1.2pt,
  left=4mm, right=4mm, top=2.5mm, bottom=2.5mm,
  drop shadow=black!10!white
]
  \centering
  {\Large\bfseries\color{primary}\faBookmark\ BÀI 66: XEM ĐỒNG HỒ. THÁNG -- NĂM}\\[1.5mm]
  {\small\itshape\color{black!75}Kĩ năng đọc đồng hồ chính xác đến từng phút, số ngày của 12 tháng trong năm và quy tắc nắm tay trên 1 trang A4}
\end{tcolorbox}

\vspace{-2mm}


\noindent
\begin{minipage}[t]{0.49\textwidth}\vspace{0pt}
\begin{tcolorbox}[
  enhanced,
  colback=blue!4!white,
  colframe=blue!75!black,
  arc=2mm, boxrule=0.8pt,
  left=2mm, right=2mm, top=1.8mm, bottom=1.8mm,
  title={\footnotesize\bfseries\color{white}\faClockO\ 1. Kĩ năng xem đồng hồ chính xác}
]
{\footnotesize
\begin{itemize}\setlength{\itemsep}{1.5pt}\setlength{\leftskip}{-2.5mm}
  \item Kim ngắn chỉ \textbf{giờ}, kim dài chỉ \textbf{phút}.
  \item Khoảng cách giữa 2 số liền nhau trên mặt đồng hồ là $\mathbf{5\text{ phút}}$.
  \item \textbf{Đọc giờ hơn:} Khi kim phút chưa vượt quá số $6$ (dưới $30$ phút).
  \item \textbf{Đọc giờ kém:} Khi kim phút đã vượt qua số $6$ (trên $30$ phút):
  $$\text{Giờ kém} = (\text{Giờ hiện tại} + 1)\text{ giờ kém } (60 - \text{số phút})\text{ phút}$$
  \item \textbf{Ví dụ:} $7$ giờ $45$ phút còn gọi là $8$ giờ kém $15$ phút.
\end{itemize}
}
\end{tcolorbox}
\end{minipage}
\hfill
\begin{minipage}[t]{0.49\textwidth}\vspace{0pt}
\begin{tcolorbox}[
  enhanced,
  colback=teal!4!white,
  colframe=teal!75!black,
  arc=2mm, boxrule=0.8pt,
  left=2mm, right=2mm, top=1.8mm, bottom=1.8mm,
  title={\footnotesize\bfseries\color{white}\faCheckSquare\ 2. Số ngày các tháng trong năm}
]
{\footnotesize
\begin{itemize}\setlength{\itemsep}{1.5pt}\setlength{\leftskip}{-2.5mm}
  \item Một năm có $\mathbf{12\text{ tháng}}$ (từ tháng 1 đến tháng 12).
  \item \textbf{Các tháng có đúng 31 ngày (7 tháng):}\\
  Tháng $1, 3, 5, 7, 8, 10, 12$.
  \item \textbf{Các tháng có đúng 30 ngày (4 tháng):}\\
  Tháng $4, 6, 9, 11$.
  \item \textbf{Tháng 2 đặc biệt:}
  \begin{itemize}\setlength{\itemsep}{0.5pt}\setlength{\leftskip}{-2.5mm}
    \item Năm thường: Tháng 2 có $\mathbf{28\text{ ngày}}$ (năm có $365$ ngày).
    \item Năm nhuận: Tháng 2 có $\mathbf{29\text{ ngày}}$ (năm có $366$ ngày).
  \end{itemize}
\end{itemize}
}
\end{tcolorbox}
\end{minipage}

\vspace{-2.5mm}
% --- KHỐI 3: SƠ ĐỒ TƯ DUY TIKZ ---
\noindent
\begin{tcolorbox}[
  enhanced,
  colback=orange!4!white,
  colframe=orange!80!black,
  arc=2mm, boxrule=0.8pt,
  left=2mm, right=2mm, top=1.5mm, bottom=1.5mm,
  title={\footnotesize\bfseries\color{white}\faSitemap\ SƠ ĐỒ TƯ DUY: MẸO NẮM TAY XÁC ĐỊNH SỐ NGÀY TRONG THÁNG}
]
\centering
\begin{tikzpicture}[
  box/.style={draw, rounded corners=2pt, align=center, font=\scriptsize\bfseries, inner sep=2.5pt, minimum height=0.75cm},
  arr/.style={-{Stealth[length=2mm]}, thick}
]
  \node[box, fill=purple!15, draw=purple!80!black, text width=2.8cm] (n1) at (0, 0) {KHỚP NỔI LỒI LÊN\\Có đúng $\mathbf{31\text{ ngày}}$\\Tháng $1, 3, 5, 7, 8, 10, 12$};
  \node[box, fill=teal!15, draw=teal!80!black, text width=2.8cm] (n2) at (4.0, 0) {RÃNH LÕM XUỐNG\\Có $\mathbf{30\text{ ngày}}$\\Tháng $4, 6, 9, 11$};
  \node[box, fill=blue!12, draw=blue!80!black, text width=2.8cm] (n3) at (8.0, 0) {RÃNH THÁNG 2\\Đặc biệt nhất\\$28$ hoặc $29$ ngày};
  \node[box, fill=orange!15, draw=orange!80!black, text width=2.4cm] (n4) at (11.2, 0) {ĐỒNG HỒ\\Kim ngắn: Giờ\\Kim dài: Phút\\Đọc hơn / kém};

  \draw[arr, purple!80!black] (n1) -- (n2);
  \draw[arr, teal!80!black] (n2) -- (n3);
  \draw[arr, blue!80!black] (n3) -- (n4);
\end{tikzpicture}
\end{tcolorbox}

\vspace{-2.5mm}
% --- KHỐI 4: VÍ DỤ & MẸO TRÁNH BẪY ---
\noindent
\begin{minipage}[t]{0.49\textwidth}\vspace{0pt}
\begin{tcolorbox}[
  enhanced,
  colback=green!4!white,
  colframe=green!60!black,
  arc=2mm, boxrule=0.8pt,
  left=2mm, right=2mm, top=1.8mm, bottom=1.8mm,
  title={\footnotesize\bfseries\color{white}\faCheckCircle\ Ví dụ thực hành chuẩn mực}
]
{\footnotesize
\textbf{1. Đọc giờ đồng hồ:}\\
Kim ngắn qua số $4$, kim dài chỉ số $8$:\\
Cách 1 (Giờ hơn): $\mathbf{4\text{ giờ } 40\text{ phút}}$.\\
Cách 2 (Giờ kém): $\mathbf{5\text{ giờ kém } 20\text{ phút}}$.\\
\textbf{2. Nhận biết số ngày:}\\
-- Tháng 7 có $\mathbf{31\text{ ngày}}$, tháng 8 có $\mathbf{31\text{ ngày}}$ (hai tháng liên tiếp cùng có 31 ngày).\\
-- Tháng 9 có $\mathbf{30\text{ ngày}}$.
}
\end{tcolorbox}
\end{minipage}
\hfill
\begin{minipage}[t]{0.49\textwidth}\vspace{0pt}
\begin{tcolorbox}[
  enhanced,
  colback=red!4!white,
  colframe=red!70!black,
  arc=2mm, boxrule=0.8pt,
  left=2mm, right=2mm, top=1.8mm, bottom=1.8mm,
  title={\footnotesize\bfseries\color{white}\faExclamationTriangle\ Mẹo tránh bẫy trắc nghiệm}
]
{\footnotesize
\begin{itemize}\setlength{\itemsep}{1.5pt}\setlength{\leftskip}{-2.5mm}
  \item \textbf{Bẫy đọc giờ kém tăng nhầm giờ:}\\
  $4$ giờ $50$ phút là $5$ giờ kém $10$ phút. Nhiều bạn đọc nhầm thành $4$ giờ kém $10$ phút là \textbf{SAI}! Giờ kém phải lấy giờ tiếp theo.
  \item \textbf{Bẫy hai tháng 7 và 8:}\\
  Hai tháng $7$ và $8$ nằm liền nhau nhưng \textbf{cả hai đều có 31 ngày} (ở khớp nổi khi đếm trên nắm tay).
\end{itemize}
}
\end{tcolorbox}
\end{minipage}


\end{document}
